\section{Immutable Parent Pointers}
\label{sec:immutable-parent-pointers}

\subsection{The ParentPointer Relation}

Central to the Recursive Spawning mechanism is the \textit{ParentPointer} relation, which establishes the unidirectional accountability link between a parent BX3 loop and each of its spawned children. We define this relation formally:

\begin{definition}[ParentPointer]
\label{def:parent-pointer}
Given a parent node $p$ and a child node $c$ produced by $p$'s spawning operation, the relation $\textit{ParentPointer}(p, c)$ is established at the moment of spawn and satisfies two governing properties:

\begin{enumerate}[noitemsep]
    \item \textbf{Uniqueness.} For any child node $c$, there exists \textit{exactly one} parent $p$ such that $\textit{ParentPointer}(p, c)$ holds.
    \item \textbf{Immutability.} Once $\textit{ParentPointer}(p, c)$ is established, it cannot be modified, reassigned, or severed for the operational lifetime of $c$.
\end{enumerate}
\end{definition}

The uniqueness property ensures that the accountability chain is unambiguous: every node in the system has exactly one authoritative parent from which it derives its Purpose and to whom it is accountable. The immutability property ensures that this chain cannot be corrupted, redirected, or broken by any subsequent operation within the spawned node, by any third-party actor, or by any network-level event.

The key architectural commitment is that the immutability of $\textit{ParentPointer}(p, c)$ is enforced by the \textbf{Fact Layer} as a physical constraint — not as a permissions check within the Bounds Engine, not as a policy rule enforced by the Purpose Layer, but as a structural property of the node's Fact Layer memory. No actor operating at a higher layer can rewrite the parent pointer after it has been set.

We denote the parent pointer of a node $a$ as $\alpha(a)$, where $\alpha(a) = p$ iff $\textit{ParentPointer}(p, a)$ holds. The function $\alpha$ is a total function on all non-root nodes and is undefined (or equivalently, returns $\bot$) for the root human node.

\subsection{The Accountability Chain}

The ParentPointer relation generates a directed acyclic graph (DAG) of accountability that we call the \textit{accountability chain}. We define this chain formally as:

\begin{definition}[Accountability Chain]
\label{def:chain}
For any node $a$ in a spawned BX3 population, the accountability chain of $a$, denoted $\textit{Chain}(a)$, is defined recursively as:

\begin{equation}
\label{eq:chain-recursion}
\textit{Chain}(a) = \{( \alpha(a),\, a )\} \;\cup\; \textit{Chain}(\alpha(a))
\end{equation}

\noindent with the base case:

\begin{equation}
\label{eq:chain-base}
\alpha(r) = \bot \quad\text{for the root human node } r \;\Rightarrow\; \textit{Chain}(r) = \emptyset
\end{equation}

where $\bot$ denotes the null parent, and the root human node is defined as the single node in the system whose parent is $\bot$.
\end{definition}

The recursion terminates at the root human node, which serves as the unique accountability anchor for the entire system. The root human has no parent, and therefore an empty chain. All other nodes have a non-empty chain whose terminal element is the root human. We prove termination:

\begin{corrolary}[Chain Termination]
For any node $a$ in a closed BX3 population with a finite spawn history, $\textit{Chain}(a)$ terminates at the root human node $r$ in a finite number of recursive applications.
\end{corrolary}

\begin{Proof}
The spawn relation is acyclic by construction: no node can spawn itself, and no node can be its own transitive ancestor. Since the spawn history is finite, the longest possible chain has depth equal to the number of spawn generations. Repeated application of the parent function therefore visits a strictly decreasing sequence of nodes in a finite set, guaranteeing termination at $\bot$. \qed
\end{Proof}

The accountability chain enables a critical inference: for any node $a$, the complete set of ancestors is recoverable by transitively applying the ParentPointer relation until the root is reached. This means that any action taken by any node in the system is traceable, through its chain of Purpose intermediaries, to a named human anchor who bears ultimate accountability.

We further define the \textit{depth} of a node as the cardinality of its accountability chain:

\begin{equation}
\label{eq:depth}
\textit{depth}(a) = |\textit{Chain}(a)|
\end{equation}

The root human node has depth $0$. In a well-formed BX3 deployment with a maximum spawn depth bound $D_{\max}$, all nodes satisfy $\textit{depth}(a) \leq D_{\max}$.

\subsection{The Root Human Anchor}

The root human node is the single human entity at the apex of the accountability hierarchy. In the context of the BX3 Framework, this is the founder, operator, or designated authority who initiates the system and retains ultimate accountability for all downstream spawns. The root human node is characterized by the following properties:

\begin{itemize}[noitemsep]
    \item \textbf{Null Parent.} $\alpha(r) = \bot$. The root human is the origin point of all authority and cannot themselves be spawned by any other node within the BX3 system.
    \item \textbf{Unconditional Accountability.} All chains terminate at the root human. The root bears final accountability for every action taken by every descendant node.
    \item \textbf{Source of Spawn Authorization.} Only the root human — or a node explicitly authorized by the root human — can authorize the creation of a new spawn event. This authorization is mediated through the Purpose Layer.
    \item \textbf{Immutable Identity.} The root human node is registered at system initialization and is cryptographically bound to the top-level Purpose Layer. It cannot be replaced, merged, or shadowed without explicit human action recorded in the Forensic Ledger.
\end{itemize}

In organizational deployments, the root human may be a founder, a board of directors, or a regulatory authority. The label does not change the structural guarantee: the root human is the unique terminus of every accountability chain, and every spawned node in the system traces back to this anchor.

The \textit{Human Root Mandate} is the architectural requirement that the root of every spawn chain must be a human-occupied node at Level 0. Without this mandate, an immutable parent pointer chain could theoretically terminate at a machine actor that was itself spawned without human authorization. The mandate closes this gap: any chain that reaches $\bot$ must have reached a human, by definition. This is enforced at the runtime level through an out-of-band human authentication ceremony for the Level 0 Human Root credential, and no machine process can manufacture this credential.

\subsection{Immutability Enforcement by the Fact Layer}

The immutability of ParentPointer is not a convention or a software policy. It is a physical constraint enforced by the Fact Layer of the node on which the ParentPointer is stored. This distinction is architecturally critical: immutability is not ``the Bounds Engine does not allow reassignment'' or ``the Purpose Layer policy prohibits it.'' It is that the Fact Layer of every node \textit{physically cannot execute} any operation that would modify, delete, or redirect an established ParentPointer.

This enforcement operates at the hardware or runtime level:

\begin{itemize}[noitemsep]
    \item \textbf{Read-only registers.} ParentPointer values are stored in Fact Layer registers or protected memory regions that are read-only at the Fact Layer boundary. No actor operating at the Bounds Engine layer can issue a write to these regions.
    \item \textbf{Separation of authorization from assignment.} The Purpose Layer can authorize a new spawn, but it cannot edit an existing ParentPointer. Authorization for spawn and immutability of existing pointers are separate, non-confusable operations performed by distinct layers.
    \item \textbf{Network-level attack resistance.} Network events, including man-in-the-middle attacks, node compromise, or replay attacks, cannot alter an established ParentPointer because the Fact Layer rejects any write operation targeting these values regardless of the requestor's credentials or network position.
    \item \textbf{Catastrophic failure resilience.} A node that experiences power loss or hardware damage does not lose its ParentPointer because the pointer is stored in non-volatile Fact Layer memory. Upon reboot, the node resumes with its chain intact.
\end{itemize}

The enforcement model follows directly from the BX3 layer separation: the Fact Layer enforces physical constraints. ParentPointer immutability is a physical constraint. Therefore, the Fact Layer enforces it. The Bounds Engine may model, reason about, or propose around the chain, but it cannot touch it.

\begin{definition}[Modification Resistance]
\label{def:modification-resistance}
A parent pointer $\alpha(c)$ is \textit{modification-resistant} if and only if all three of the following hold:

\begin{enumerate}[noitemsep]
    \item Any write to $\alpha(c)$ after initialization requires Fact Layer authorization with a corresponding Purpose Layer approval and a Forensic Ledger entry.
    \item The write is atomic with respect to the ledger update; the pointer is not updated unless the ledger entry is confirmed.
    \item Any unauthorized write attempt is rejected at the Fact Layer boundary, logged as a \texttt{PTR\_WRITE\_DENIED} event in the Forensic Ledger, and triggers the Bailout Protocol.
\end{enumerate}
\end{definition}

The atomicity condition in item (2) is critical: it prevents a scenario in which the pointer is updated but the ledger write fails, creating an inconsistent state where the Fact Layer and Forensic Ledger diverge. The atomicity guarantee means that a ParentPointer change is impossible without a corresponding attestation, and a ledger entry is impossible without a confirmed pointer state.

\subsection{Spawn Authorization via the Purpose Layer}

While existing ParentPointers are immutable, the creation of a \textit{new} ParentPointer — i.e., a new spawn event — requires explicit authorization from the Purpose Layer. This is the sole mechanism by which the accountability chain grows. The Purpose Layer mediates spawn authorization through the following sequence:

\begin{enumerate}[noitemsep]
    \item \textbf{Request.} A parent node's Bounds Engine determines that a new child node is required to handle a specific local context or workload. It formulates a spawn request addressed to its own Purpose Layer, specifying the authorized scope and constraint boundaries for the proposed child.
    \item \textbf{Authorization Check.} The Purpose Layer evaluates the request against its registered mission parameters, constraint envelope, and resource allocation policy. It determines whether the parent is authorized to spawn at all, and if so, what scope of Purpose the child will carry. If the request falls outside the parent's authorized spawn mandate, the Purpose Layer denies the request and logs a \texttt{SPAWN\_DENIED} event.
    \item \textbf{Worksheet Generation.} Upon authorization, the parent Bounds Engine generates a \textit{Worksheet} — a containerized, self-contained logic set that encapsulates the authorized Purpose for the specific local context. The Worksheet includes the immutable ParentPointer value and a cryptographically signed assertion of authorization provenance, including the authorizing Purpose Layer's identity and the timestamp of authorization.
    \item \textbf{Deployment.} The parent deploys the Worksheet to the child node over-the-air (OTA). The child node instantiates its BX3 loop using the received Worksheet, establishing its Fact Layer with the authorized constraint set and the immutable ParentPointer pre-loaded. The child is instantiated with its $\alpha$ value already set — it does not choose its parent.
    \item \textbf{Chain Extension.} The new child node is now live with $\textit{ParentPointer}(\text{parent}, \text{child})$ established and immutable from the moment of instantiation. The Fact Layer records a \texttt{CHAIN\_EXTEND} event in the Forensic Ledger.
\end{enumerate}

The Purpose Layer's role in spawn authorization ensures that the chain grows in a controlled, auditable manner. Every spawn event is logged in the Forensic Ledger with the authorizing Purpose Layer reference, the authorized scope, and the timestamp. No chain extension occurs without a Purpose Layer authorization record. This prevents unauthorized spawn proliferation — a failure mode documented as \textit{agent sprawl} in multi-agent systems literature \cite{moore2025hmas}.

This separation of concerns is critical: the Purpose Layer authorizes, but the Fact Layer assigns. The parent does not choose its child — it requests authorization to produce a child, and the Fact Layer performs the atomic assignment of the parent pointer at the moment of child instantiation. This prevents a parent from claiming to have spawned a node whose parent pointer was set to a different actor, and it prevents a child from falsifying its own ancestry.

\begin{property}[Authorization-Provenance Binding]
Every child node $c$ carries a cryptographically signed provenance record asserting that $\textit{ParentPointer}(p, c)$ was established under the authorization of $p$'s Purpose Layer. This record is embedded in the child's Worksheet and verified by the Fact Layer at instantiation. The provenance binding ensures that the parent pointer cannot be disputed: the child was instantiated with that pointer, and the pointer was set by the parent under Purpose Layer authorization.
\end{property}

\subsection{Chain Traversal Algorithm}

The accountability chain must be traversable by any authorized actor in the system, including auditors, regulators, and the Purpose Layer itself. We define a canonical chain traversal algorithm that operates with guaranteed termination at the root human:

\begin{algorithm}[htbp]
\caption{Accountability Chain Traversal}
\label{alg:chain-traversal}
\begin{enumerate}[noitemsep]
    \item \textbf{Input:} A node identifier $a$ for which the full accountability chain is required.
    \item \textbf{Output:} An ordered list $[r,\, p_1,\, p_2,\, \dots,\, p_n,\, a]$ where $r$ is the root human node, $p_i$ are intermediate parent nodes, and $a$ is the target node.
    \item \textbf{Procedure:}
    \begin{enumerate}[noitemsep]
        \item Initialize an empty list $\texttt{chain}$.
        \item Set $\texttt{current} \gets a$.
        \item \textbf{while} $\texttt{current} \neq \bot$ \textbf{do}:
        \begin{enumerate}[noitemsep]
            \item $\texttt{append}(\texttt{chain}, \texttt{current})$.
            \item $\texttt{parent} \gets \textit{ReadParent}(\texttt{current})$ \hfill \textit{// Fact Layer read operation.}
            \item \textbf{if} $\texttt{parent} = \bot$ \textbf{then}:
            \begin{enumerate}[noitemsep]
                \item \textbf{assert} $\texttt{current}$ is the root human $r$.
                \item \textbf{break}.
            \end{enumerate}
            \item \textbf{else}: $\texttt{current} \gets \texttt{parent}$.
        \endenumerate}
        \item \textbf{return} $\texttt{reverse}(\texttt{chain})$ \hfill \textit{// root first, target last.}
    \endenumerate}
\end{enumerate}
\end{algorithm}

The algorithm terminates because the ParentPointer relation is acyclic: a node cannot be its own ancestor, and the chain is finite by construction (the maximum depth is the number of spawn generations). The Fact Layer enforces that $\textit{ReadParent}()$ always returns the stored pointer value and never a computed or modified value, preventing any manipulation of the traversal from the Bounds Engine layer.

The reverse-ordered output $[r, p_1, p_2, \dots, p_n, a]$ is the standard representation used in Forensic Ledger entries, audit reports, and escalation notifications: root human first, immediate parent second, and the node in question last. This ordering makes the chain immediately human-readable and supports the upstream accountability guarantee that every node is traceable to its human root.

\begin{property}[Escalation Time Bound]
For any spawned population with a fixed maximum spawn depth $D_{\max}$, the chain traversal algorithm completes in $O(D_{\max})$ steps. In typical BX3 deployments, $D_{\max} \leq 5$, meaning the traversal requires at most five pointer dereferences to reach a human.
\end{property}

The constant-time escalation property is a direct consequence of immutable parent pointers combined with depth-bounded spawning. Mutable parent pointers would allow the chain to be extended retroactively, making the effective depth potentially unbounded and the worst-case escalation time unbounded as well. Immutability is what makes the depth bound a structural guarantee, not merely a policy.

This bound is particularly significant for the Bailout Protocol: when a node encounters an unresolvable state, it must escalate to its parent. With $O(D_{\max})$ traversal, the maximum number of escalation hops is bounded, ensuring that unresolved states reach a human within a predictable time window regardless of population size.

\subsection{Forensic Ledger Integration}

Every ParentPointer read during chain traversal is logged in the Forensic Ledger as a \texttt{CHAIN\_TRAVERSE} event. Every spawn event that creates a new ParentPointer is logged as a \texttt{CHAIN\_EXTEND} event. These events are recorded simultaneously on the Purpose, Bounds Engine, and Fact planes, providing three-plane attestation of chain state. If any plane's log diverges from the others, a \texttt{CHAIN\_INTEGRITY\_FAILURE} flag is raised and a Bailout Protocol signal is sent to the nearest Purpose Layer ancestor.

The three-plane attestation requirement is non-negotiable: a chain integrity claim requires agreement across all three planes. A \texttt{CHAIN\_EXTEND} event recorded on the Fact plane but absent from the Purpose plane indicates a potential compromise of the Purpose Layer's authorization logging. A chain traversal record on the Bounds Engine plane but absent from the Fact plane indicates a potential attempt to spoof the chain through a compromised Bounds Engine. Divergence at any plane triggers automatic escalation and a potential Bailout Protocol activation.

This multi-plane attestation model draws on established principles in distributed systems integrity monitoring. The simultaneous logging across planes ensures that no single plane can be compromised to falsify chain state without detection by the remaining two.

\subsection{Design Rationale}

The immutable parent-pointer architecture draws on a foundational principle from cybernetics: a control system without a feedback loop cannot distinguish between successful execution and failure \cite{wiener1948}. The parent pointer is the feedback loop in the accountability architecture — the mechanism by which the system continuously verifies that every active agent remains connected to a named human decision.

The physical analogy is genealogical lineage: every individual has biological parents whose assignment is immutable once established, and this lineage cannot be retroactively rewritten by the individual or any third party. The immutable parent pointer serves the same structural function as biological parentage in genealogical accountability: it is the permanent, non-negotiable anchor that connects any agent to its origin. Just as no individual can choose their biological parents, no agent can choose or alter its parent pointer after spawn.

The design rationale for immutability enforcement at the Fact Layer reflects the framework's core architectural principle: accountability must be enforced at the layer that cannot reason and therefore cannot be deceived by its own output. The Fact Layer executes; it does not deliberate. When it records that $\textit{ParentPointer}(p, c)$ has been established, that record is a physical fact about the state of the system — not a claim that could be falsified by subsequent reasoning. This is the same principle that makes the Forensic Ledger tamper-evident: the layer that records facts cannot be tricked by the layer that reasons about them.

The chain traversal algorithm provides the operational counterpart to the structural immutability guarantee. Immutability ensures that parent-pointer records cannot be altered after creation; chain traversal ensures that the complete, unbroken chain from any agent to its Human Root can always be reconstructed on demand. Together, these mechanisms make the accountability chain not merely a policy aspiration but a verifiable, machine-auditable property of the runtime — one that can be independently verified by external auditors, regulators, or dispute-resolution processes without relying on any agent to self-report its own chain.

The immutable parent-pointer architecture is what makes recursive spawning in BX3 fundamentally different from agent-forking in conventional multi-agent systems. In conventional systems, a spawned agent may retain a reference to its parent but can re-parent itself, spawn unauthorized children, or sever the link under certain runtime conditions \cite{agentspawn2025,firmhive2025}. In BX3, the parent pointer is architecturally immutable once assigned — it is a permanent structural property of the agent, not a mutable reference that can be updated by subsequent operations. This distinction transforms the spawn chain from a flexible delegation heuristic into a verifiable accountability infrastructure, and it is the basis for the downstream guarantees that the Bailout Protocol and the Training Wheels Protocol depend upon.